\chapter{Conclusion and Future Enhancements}

\section{Conclusion}

\project{} set out to remove the dedicated identity hardware from a parking gate and
replace it with a phone that the driver already owns. The delivered system does that
without weakening the access decision: the phone carries a session identifier, the
gate asks the backend, and the backend alone decides. Entry starts the billing clock,
payment freezes the amount, and the exit verdict checks the payment rather than the
clock. A camera, an RFID reader, a ticket dispenser and a vehicle sensor are all
absent, and the entire site is one ESP32 and one servo.

The work produced four artefacts that can be evaluated separately. The backend is a
layered Spring Boot application with three credential schemes, a transactional
session lifecycle, per-second billing and a device-command queue. The mobile
application is a typed Expo client that owns no business rules and displays the
backend's own messages. The firmware is an Arduino sketch whose transports and
actuator sit behind interfaces, so the same gate logic runs over Wi-Fi, over a serial
demonstration bridge, or against an LED stand-in. The verification workflow is a
Compose profile that boots the real image as a readiness gate, runs 118 tests under a
100\% service-line coverage gate, and publishes JUnit and JaCoCo reports even when the
run fails.

Three conclusions are worth carrying forward. First, \emph{the correct place for the
access decision is the server}, and moving it there makes the edge dramatically
simpler rather than more complex. Second, \emph{duration and payment are the only
inputs the exit gate needs}; once that is true, a single gate can serve both
directions and the facility needs no second mechanism. Third, \emph{reproducible
verification is a feature}: the container workflow turned every later change into a
checkable one and made the coverage figure a constraint rather than a report.

The system is complete for its stated scope. What remains is not unfinished work but
the next layer of hardening: real payments, transport security, richer gate hardware
and the operational surfaces that a commercial deployment would require.

\section{Future Enhancements}

\begin{table}[htbp]
  \centering
  \small
  \begin{tabularx}{\textwidth}{@{}lXl@{}}
    \toprule
    \textbf{Enhancement} & \textbf{Description and why it matters} & \textbf{Priority} \\
    \midrule
    Real payment gateway & Replace the mock provider with an implementation of
      \code{PaymentService} that supports authorisation, failure and refunds. The
      session lifecycle already models the states a gateway needs. & High \\
    Transport security & Serve the API over TLS and move the gate to HTTPS; consider
      BLE pairing with an authenticated session token instead of a plaintext
      identifier. & High \\
    Hardware validation & Execute and record the servo, watchdog and BLE bench
      procedures on a physical gate, and add an automated end-to-end test with a
      development board on the bench. & High \\
    Escalating enforcement & Add sensors, a beam or a loop detector so the gate can
      detect a tailgating vehicle and refuse to close on an obstruction. & Medium \\
    Multiple gates and lanes & Generalise the device model to entry and exit gates
      per facility and to more than one servo, with a lane identifier in the
      verdict. & Medium \\
    Operator console & A small web surface for facilities, pricing, device rotation
      and live gate state, replacing curl-style administration. & Medium \\
    Device secret rotation & A rotation endpoint that returns a new secret once,
      invalidates the old hash immediately and records the event. & Medium \\
    Over-the-air updates & Signed firmware updates for the gate, with rollback, so a
      provisioned site does not need physical access to change behaviour. & Low \\
    Occupancy and analytics & Optional bay sensors or camera counts to expose
      utilisation and revenue trends without making access depend on them. & Low \\
    Background BLE and wallet cards & iOS background operation and a fallback NFC or
      card credential for drivers without the application. & Low \\
    \bottomrule
  \end{tabularx}
  \caption{Enhancement roadmap, ordered by the risk each item removes.}
  \label{tab:future}
\end{table}

The ordering follows one rule: remove the assumptions that would be unacceptable in
production before adding capability. A real payment provider and transport security
address the two places where the current system is explicitly a prototype. Hardware
validation converts documented procedures into measured evidence. Everything else
expands the product; none of it changes the architecture, because the decision that
matters is already in the right place.
